\section{Interpretable Models}
Using inherently interpretable algorithms.
\subsection{Linear Regression}
Predicts target as weighted sum of inputs.
\begin{equation}
	y = \beta_0 + \beta_1 x_1 + \dots + \beta_p x_p + \epsilon
	\label{eq:linear_regression}
\end{equation}
Estimated usually by Ordinary Least Squares (OLS)
\begin{equation}
	\hat \beta = \argmin{\beta_0, \dots , \beta_p} \sum_{i=1}^n \lr{y^{(i)} - \lr{\beta_0 + \sum_{j=1}^p}\beta_jx_j^{(i)}}²
	\label{eq:OLS}
\end{equation}
Each weight represents the effect of one feature, holding others constant.
We can compute \textbf{confidence intervals} for each coefficient.
\paragraph{Regression Output}
\begin{table}[H]
	\caption{Regression Output}\label{tab:regression_output}
	\begin{center}
		\begin{tabular}{l|r|r|r}
			\hline
			Feature     & Weight$\hat \beta_j$ & SE  & $|t|$ \\
			\hline
			Temperature & 110.7                & 7.0 & 15.7  \\
			Humidity    & -17.4                & 3.2 & 5.5   \\
			\hline
		\end{tabular}
	\end{center}
\end{table}
\begin{itemize}
	\item Weight: Estimated influence of $x_i$ over the outcome $y$
	\item SE (Standard Error): Uncertainty (variance) of the estimated coefficient
	\item $|t|$ t-statistic: Strength of evidence that $\beta_j \neq 0$
\end{itemize}
\begin{equation}
	t_{\beta_j} = \frac{\hat \beta_j}{SE(\hat \beta_j)}
	\label{eq:t-statistic}
\end{equation}
Large, certain weights = important features. Small or uncertain weights = weak evidence. Features wit highest $|t|$ values are most significant.

\begin{itemize}
	\item Effect Plot: $effect_j^i = \beta_j \cdot x_j^i$ Contribution of feature j to prediction for instance i.
	\item Different encodings for variables.
\end{itemize}

\subsection{Shrinkage Models}
Using models that keep only the relevant features helps with interpretability.

\subsubsection{Lasso - Least Absolute Shrinkage and Selection Operator}
\begin{itemize}
	\item Adds L1 Norm, Performs feature selection, regularization
	\item Enforces sparsity in Linear model
\end{itemize}
\begin{equation}
	\underset{\beta}{min} \frac{1}{n}\sum:{i=1}^n\lr{y^{(i)}-x_i^T \beta }² + \lambda \norm{\beta}_1
	\label{eq:LASSO-L1}
\end{equation}
As $\lambda$ gets increased, only the most important features remain. This increases interpretability.

\paragraph{About Linear Regression}
\begin{itemize}
	\item Perfectly understood
	\item Only captures linear relationships
	\item Ignores feature interactions, must be added manually ($x_1x_2, x²$ etc.)
	\item Lower accuracy compared to more flexible Models.
\end{itemize}


\subsection{Logistic Regression}
\begin{itemize}
	\item Classification problems with two classes.
	\item Models Probability of one class given the features.
	\item Extension of linear regresion for classification.
	\item Uses the output of a linear regression and maps it to a Probability.
\end{itemize}
\paragraph{Logistic / Sigmoid function}
\begin{equation}
	logistic\lr{\eta} = \frac{1}{1+ \exp{-\eta}}
	\label{eq:Log-sigmoid-fun}
\end{equation}

\paragraph{Logistic Model}
\begin{equation}
	P\lr{y^{(i)} = 1}  = \frac{1}{1 + \exp (- (\beta_0 + \beta_1 x_1^{(i)} + \dots + \beta_p x_p^{(i)}))}
	\label{eq:logisticregression}
\end{equation}
\begin{itemize}
	\item Linear predictor is transformed via the logistic function.
	\item Coeffictions do not act linearly on the probability. Interpretation is happens on the log-odds scale.
	\item A one unit increase in $x_j$ multiplies the odds of $y=1$ vs. $y=0$ by a factor of $\exp(\beta_j)$
	\item Predicts probabilities, which alows for better decision making. Pred with 99\% probability is more confident than one with 51\%.
	\item Relationships between features and outcomes remain transparent.
\end{itemize}


\subsection{Generalized Linear Models}
Linear Regression





\subsection{Generalized Additive Models}
